\documentclass[11pt]{article}
\usepackage[margin=1in]{geometry}
\usepackage{listings}
\usepackage{xcolor}
\usepackage{graphicx}
\usepackage{float}

\lstset{
    language=C++,
    basicstyle=\ttfamily\small,
    keywordstyle=\color{blue},
    commentstyle=\color{gray},
    stringstyle=\color{red},
    breaklines=true,
    showstringspaces=false,
    tabsize=4,
    numbers=left,
    numberstyle=\tiny,
    frame=tb
}


\title{Laborator 3 - SO}
\author{Rus Lucian-Florin}
\date{April 22, 2026}

\begin{document}

\maketitle

\newpage
\section*{Comenzi laborator}

\begin{figure}[H]
    \centering
    \includegraphics[width=0.8\textwidth]{1.png}
    \caption{Set comenzi laborator 1}
    \label{fig:1}
\end{figure}

\begin{figure}[H]
    \centering
    \includegraphics[width=0.8\textwidth]{2.png}
    \caption{Set comenzi laborator 2}
    \label{fig:2}
\end{figure}

\newpage
\section*{Teme}

\subsection*{Problema 2}

\begin{lstlisting}
#include <dirent.h>
#include <stdio.h>
#include <sys/stat.h>
#include <sys/types.h>
#include <unistd.h>

int main(void) {
    char curr[256];
    // ideal era sa folosim `get_current_dir_name`, dar folosim `getcwd` ca sa ramanem posix-compliant
    if (getcwd(curr, 256) == NULL) {
        printf("eroare la deschidere\n");
        return -1;
    }
    DIR *dir = opendir(curr);

    struct dirent *entry;
    // entry e un pointer la dir-ul curent, daca reapelam functia, inaintam pointer-ul la toate celelalte entitati care se afla in
    // directorul target
    while ((entry = readdir(dir)) != NULL) {
        struct stat buf;
        if (lstat(entry->d_name, &buf) < 0) {
            printf("lstat error for %s", entry->d_name);
            continue;
        }

        // facem stat ca sa verificam bit-ul de executie. daca e setat pentru oricare dintre situatiile de mai jos printam textul
        // cu un format verde bold, asemanator unui terminal clasic, dupa care resetam culoarea textului
        if (buf.st_mode & (S_IXGRP | S_IXOTH | S_IXUSR)) {
            printf("\033[1;32m%s\033[0m\n", entry->d_name);
        }
        else {
            printf("%s\n", entry->d_name);
        }
    }
    closedir(dir);

    return 0;
}
\end{lstlisting}

\newpage
Snippet-ul de cod de pe pagina anterioara contine implementarea comenzii \textbf{nls}, reprezentand o implementare proprie a comenzii standard \textbf{ls} din Linux. Output-ul ambelor comenzi se poate vedea in imaginea alaturata. Pentru o simulare cat mai apropiata a comportamentului, s-au folosit si coduri de culoare.

\begin{figure}[H]
    \centering
    \includegraphics[width=0.8\textwidth]{3.png}
    \caption{Output comenzi \textbf{ls} si \textbf{nls}}
    \label{fig:3}
\end{figure}

\newpage

\subsection*{Problema 3}

\begin{lstlisting}
static int myftw(char *pathname, Myfunc *func) {

    // fullpath = path_alloc(NULL);	/* malloc's for PATH_MAX+1 bytes */
    // strcpy(fullpath, pathname);	/* initialize fullpath */

    // pornim de la sursa
    if (chdir(pathname) < 0) {
        err_ret("error chdir %s", pathname);
        return -1;
    }
    return dopath(func);
}
static int dopath(Myfunc *func) {
    struct stat    statbuf;
    struct dirent *dirp;
    DIR           *dp;
    int            ret = 0;
    char          *ptr;

    // luam directorul curent
    char curr[256] = {0};
    if (getcwd(curr, 256) == NULL) {
        printf("eroare la deschidere\n");
        return func(curr, &statbuf, FTW_NS);
    }

    // printf("current: %s\n", curr);

    if (lstat(curr, &statbuf) < 0)
        return func(curr, &statbuf, FTW_NS); /* stat error */

    if (S_ISDIR(statbuf.st_mode) == 0)
        return func(curr, &statbuf, FTW_F); /* not a directory */

    /*
     * It's a directory.  First call func() for the directory,
     * then process each filename in the directory.
     */

    if ((ret = func(curr, &statbuf, FTW_D)) != 0)
        return ret;

    // ptr = fullpath + strlen(fullpath);	/* point to end of fullpath */
    // *ptr++ = '/';
    // *ptr = 0;

    if ((dp = opendir(curr)) == NULL)
        return func(curr, &statbuf, FTW_DNR);

    while ((dirp = readdir(dp)) != NULL) {
        if (strcmp(dirp->d_name, ".") == 0 || strcmp(dirp->d_name, "..") == 0)
            continue; /* ignore dot and dot-dot */

        // strcpy(ptr, dirp->d_name);	/* append name after slash */

        struct stat entry;
        if (lstat(dirp->d_name, &entry) < 0)
            continue;

        // cum folosim chdir, trebuie sa ne asiguram ca ne intoarcem corect
        if (S_ISDIR(entry.st_mode)) {
            if (chdir(dirp->d_name) == 0) {
                if ((ret = dopath(func)) != 0) { /* recursive */
                    chdir("..");
                    break; /* time to leave */
                }
                chdir("..");
            }
        }
        else {
            if ((ret = func(dirp->d_name, &entry, FTW_F)) != 0)
                break;
        }
    }

    // ptr[-1] = 0;	/* erase everything from slash onwards */

    if (closedir(dp) < 0)
        err_ret("can't close directory %s", curr);

    return ret;
}
\end{lstlisting}

\begin{figure}[H]
    \centering
    \includegraphics[width=0.8\textwidth]{4.png}
    \caption{Output comenzi \textbf{ftw4} si \textbf{ftw4-2}}
    \label{fig:4}
\end{figure}

Din imagine, se poate observa ca implementarea \textbf{ftw4-2}, implementarea care contine comanda de \textit{chdir} este putin mai rapida. Probabil, scalat la un director mai mare, timpii de executie ai comenzii implementate pentru tema ar fi simtitor mai buni. 
\end{document}


